\section{Polymorphic arrays: adapting the integer library}
\label{sec:poly-arrays}

The integer array of Section~\ref{sec:arrays} remains in place, with its
original code, proofs, and clients. This separate example, in
\path{example/poly_array/rocq/poly_array.v}, adapts its representation and
proofs to arbitrary values. The interface is familiar:
\begin{lstlisting}
class PolyArray<T> {
    static func make(n: Int64, seed: T): PolyArray<T>
    func length(): Int64
    func get(i: Int64): T
    func set(i: Int64, x: T): Unit
}
\end{lstlisting}
These are signatures; the following listing is a runnable use of the
generated methods. The type argument belongs to the class; it is written
on the factory owner and inferred from the receiver for instance calls.
\begin{lstlisting}
let flags = PolyArray<Bool>.make(3, false)
flags.set(1, true)
let first = flags.get(0)  // false
let middle = flags.get(1) // true
let count = flags.length() // 3
\end{lstlisting}
The initial value is explicit because an arbitrary type has no universal
zero. For integers, a zero seed matches the old initial contents. For booleans, a seed of \texttt{false} gives three false
entries. The verified constructor requires $n>0$, just as the old core
allocation does, and each verified access requires $0\leq i<n$.
The native implementation may support more inputs; those extra behaviours
are not claims of these contracts.

\subsection{The same representation, with values instead of integers}

Write $\mathsf{PA}(xs,a)$ for \texttt{is\_poly\_array xs a}. Its shape is
the old array predicate:
\[
\begin{aligned}
\mathsf{cells}(b,xs) &\;=\;
  \bigsep_{i\mapsto v\in xs} (b+i)\ptsr v,\\
\mathsf{PA}(xs,a) &\;=\;\exists o,b.\;
  \pure{a=\mathsf{obj}(o)}\sep
  o[0]\ptso |xs|\sep o[1]\ptso\mathsf{ptr}(b)\\
 &\hspace{6em}\sep b\ptsb |xs|\sep\mathsf{cells}(b,xs).
\end{aligned}
\]
The mathematical list now has type \texttt{list val}; the old list had
type \texttt{list Z}. A cell can contain an integer, a Boolean, a struct,
an object reference, or a closed function value. The theorem
\texttt{poly\_int\_adaptation} makes the relationship precise:
\[
\mathsf{PA}(\mathsf{map}\;\mathsf{LitInt}\;zs,a)
\quad\dashv\vdash\quad \mathsf{isArray}(zs,a).
\]
Here $\mathsf{LitInt}$ abbreviates the embedding of an integer into a CJR
value. The theorem is an equivalence of ownership predicates in the core
model. It is not an assertion that the two native storage layouts coincide.

The type argument $T$ is absent from this predicate. It remains in the
generic source and in printed Cangjie, while the proof reasons uniformly
about $xs$. For example, the proof of a load needs ownership of cell $i$;
it need not know whether the cell contains $42$ or \texttt{true}.

\subsection{Initialization and erasure}

The source type \texttt{GTBuffer t} and expressions
\texttt{GEBufferMake}, \texttt{GEBufferGet}, and \texttt{GEBufferSet}
describe initialized indexed storage. Their type annotations undergo the
same substitution and scope checks as other generic types. Erasure removes
the element type. A read becomes \texttt{Load (Offset b i)}, a write
becomes \texttt{Store (Offset b i) v}, and construction calls a closed
helper defined in \path{theories/buffer.v} using existing core operations.
There is no additional reduction rule or logical axiom.

The helper allocates $n$ zero words, then fills every cell with the seed.
Its recursive fill visits $n-1,n-2,\ldots,0$. To prove the successor case,
split off the last zero cell, write the seed there, apply the induction
hypothesis to the prefix, and combine the two pieces. Thus
\texttt{wp\_buffer\_fill} changes
$\mathsf{cells}(b,\mathsf{replicate}(n,0))$ into
$\mathsf{cells}(b,\mathsf{replicate}(n,v))$.
The constructor then installs the length and pointer in a new object,
as the monomorphic constructor did.

The seed and values supplied to \texttt{set} must be closed under
term substitution. In Rocq this is
\texttt{closed\_val v := forall x w, subst\_val x w v = v}.
Literals and object references satisfy it immediately. For functions, it
rules out an unresolved free local variable inside the stored function.
The condition matters because the existing core substitutes through values
while unpacking call arguments; it is unrelated to duplicating ownership.

The factory is a \texttt{GStaticFunction} owned by \texttt{PolyArray},
called with \texttt{GEStaticCall}. It has no receiver. The three instance
methods use \texttt{GEMethod}; their explicit first source parameter is
printed as \texttt{this}. Class parameters are in scope in all four bodies.
Erasure resolves qualified names such as \texttt{PolyArray.make} and
\texttt{PolyArray.get}, keeping the same argument order and runtime bodies.

Both the type-instantiation erasure theorem and the term-substitution
erasure theorem cover the new storage forms. Consequently
$\GWpre(e,\Phi)$ still means $\Wpre(\mathsf{erase}(e),\Phi)$,
and the existing adequacy theorem still supplies safety. The concrete
Boolean-array constructor is covered by \texttt{poly\_main\_adequate}.

\subsection{The four contracts and their proofs}

Suppressing the class type argument, writing an instance receiver as the
first argument, and writing $xs[i\leftarrow x]$
for list update, the structural contracts are:
\[
\begin{aligned}
\mathsf{emp} &\vdash \GWpre(\mathsf{make}(n,v),
  \lambda a.\mathsf{PA}(\mathsf{replicate}(n,v),a)),\\
\mathsf{PA}(xs,a) &\vdash \GWpre(\mathsf{length}(a),
  \lambda r.\pure{r=|xs|}\sep\mathsf{PA}(xs,a)),\\
\mathsf{PA}(xs,a) &\vdash \GWpre(\mathsf{get}(a,i),
  \lambda r.\pure{r=xs_i}\sep\mathsf{PA}(xs,a)),\\
\mathsf{PA}(xs,a) &\vdash \GWpre(\mathsf{set}(a,i,x),
  \lambda r.\pure{r=()}\sep\mathsf{PA}(xs[i\leftarrow x],a)).
\end{aligned}
\]
The first line assumes $n>0$ and a closed seed; the final two assume
$xs!!i=\mathsf{Some}(xs_i)$; the last also assumes a closed new value.
The statements appear as \texttt{poly\_make\_spec},
\texttt{poly\_length\_spec}, \texttt{poly\_get\_spec}, and
\texttt{poly\_set\_spec}.

The get and set proofs closely follow the integer proofs. Unpack the
arguments, read the data field, calculate the offset, and use the list
lookup to obtain exactly one points-to assertion. A load returns that
assertion unchanged; a store returns it with the new value. The closing
continuation supplied by the big-separating-conjunction lemma rebuilds
the array. Since list update preserves length, the length field and block
metadata remain valid.

The derived results recover the usual array equations:
\texttt{poly\_set\_get\_same} reads back the new value;
\texttt{poly\_set\_get\_other} preserves a different index;
\texttt{poly\_set\_length} preserves length; and
\texttt{poly\_fresh\_get} supplies the seed at every valid fresh index.
For the Boolean listing above, the list changes from
$[\False,\False,\False]$ to $[\False,\True,\False]$.
No integer-specific arithmetic on the elements enters these proofs.

\subsection{When elements themselves own resources}

Owning an array cell containing an object reference does not by itself
own the referenced object's fields. To express such ownership, choose an
element representation $R:\mathsf{val}\to\mathsf{iProp}$ and define
\[
  \mathsf{PA}_R(xs,a)
  =\mathsf{PA}(xs,a)\sep\bigsep_{v\in xs} R(v).
\]
This is \texttt{poly\_array\_rep}. For booleans, $R$ might be a pure fact
that the value is a Boolean. For an array of boxes, $R(v)$ might own the
box field. For an array of arrays, it can be another array predicate.

The ownership-aware get contract returns the selected resource together
with a continuation that restores the whole array:
\[
 \mathsf{PA}_R(xs,a)\vdash
 \GWpre(\mathsf{get}(a,i),
 \lambda v.\pure{v=xs_i}\sep R(v)\sep
       (R(v)\wand\mathsf{PA}_R(xs,a))).
\]
Think of opening an envelope: the selected resource is available, but the
envelope cannot be closed until that resource is returned. There is no
second copy of $R(v)$ left inside a complete $\mathsf{PA}_R$.
This is \texttt{poly\_get\_owned\_spec}.

An ownership-aware update consumes the new element resource and returns
the old one. With $xs_i=y$ and closed $x$,
\[
\begin{aligned}
 \mathsf{PA}_R(xs,a)\sep R(x)\vdash
 \GWpre(\mathsf{set}(a,i,x),\lambda r.\;&\pure{r=()}\sep
 \mathsf{PA}_R(xs[i\leftarrow x],a)\sep R(y)).
\end{aligned}
\]
The returned $R(y)$ is a logical resource, not an extra runtime return
value. The caller already knows $y$ from the precondition and can retain,
reuse, or dispose of its resource according to its own specification.

Construction requires special care. Repeating a seed $v$ creates $n$
copies of its value, so \texttt{poly\_make\_owned\_spec} requires
$\bigsep_{k<n} R(v)$. A single $R(v)$ suffices only when it is persistent;
\texttt{poly\_make\_persistent\_spec} states that case explicitly.
For an exclusively owned box, two copies of its field ownership are
unavailable. The program can still construct an array of two aliases,
but its appropriate contract keeps the box ownership once outside the
array's structural predicate.

\begin{lstlisting}
let inner = PolyArray<Bool>.make(2, false)
let outer = PolyArray<PolyArray<Bool>>.make(2, inner)
outer.get(0).set(1, true)
// Both entries contain the same inner reference:
let seen = outer.get(1).get(1)
// seen == true
\end{lstlisting}
The outer cells are distinct; the referenced inner array is shared.
For a singleton outer array, \texttt{poly\_nested\_singleton\_spec}
instead transfers the one inner-array resource into
$\mathsf{PA}_{\mathsf{PA}(xs,-)}([a],b)$ without any duplication.

\subsection{Native storage and validation}

The printer emits \texttt{GTBuffer T} as \texttt{CjrBuffer<T>}, an alias
for \texttt{std.core.Array<T>}. This supports Boolean, struct, function,
and reference elements without inventing a native byte width for $T$.
Cangjie 1.0.5 rejects an unconstrained \texttt{CPointer<T>}: its element
must satisfy the native \texttt{CType} requirement. The new example uses
managed storage; the old integer example retains its raw storage.

The core's one word per value is an abstract model of indexed storage,
not a proof of the native array implementation. The correspondence between
the derived allocation-and-fill helper and the managed constructor remains
part of the printer and native-runtime trust boundary. Verified lengths
must fit native \texttt{Int64}, indexes must satisfy the contracts, and
allocation failure and garbage collection are not modeled. The native
array's bounds checks remain enabled.

Buffer receiver annotations are checked by typed helper calls. A store
evaluates its buffer, index, and value as three arguments before performing
the assignment, retaining the core's evaluation order. An executable
side-effect trace checks all seven operands across construction, write,
and read, and obtains \texttt{1234567}. Tests cover wrong
element annotations and incompatible stores as well as successful
instantiations. Eight native cases exercise seeded construction, length,
writes at both ends, independent arrays, Boolean and integer contents,
nested arrays and aliasing, generic structs, functions, and unit values.
Rocq regressions check nested substitution, scope and arity rejection,
erasure laws, integer and Boolean round trips, and nested ownership.

\subsection{Comparison with the monomorphic array}
\label{sec:poly-array-benchmark}

The paired benchmark instantiates the new array at \texttt{Int64} and
compiles it with the unchanged generated integer array. It removes only
their demonstration entry points to link the two libraries in one test
executable. Source hashes include both complete generated files and that
small build transformation.

The workloads and sizes match the existing array benchmark: permuted
indexed reads and permuted writes followed by sequential reads, at
$n=256,4096,65536$. Setup populates both arrays before timing. Each process
warms both implementations and calibrates batches towards 20 milliseconds;
seven samples alternate their execution order. Three independent processes
give 21 samples per implementation at each workload and size, or 252
measurements. Every measured batch checks its checksum. Compilation,
allocation, initialization, and raw-block cleanup are outside the timer.
The timer is \texttt{std.time.MonoTime}; there is no core pinning or forced
garbage collection. Medians and interquartile ranges pool the 21 samples
for each configuration. In the write/read workload, one operation means
one write or one read, so a complete pass counts as $2n$ operations.

The 29 October 2026 run used Cangjie 1.0.5 with \texttt{-O2}. The
table shows median nanoseconds per operation, with the interquartile
range in parentheses. A Poly/Mono ratio above one means that the
polymorphic implementation took longer in this run.
\begin{center}
\small
\begin{tabular}{lrrrr}
\hline
Workload & $n$ & Mono ns/op & Poly ns/op & Poly/Mono \\
\hline
Read & 256 & 1.01 (0.01) & 0.67 (0.00) & 0.668 \\
Read & 4,096 & 0.98 (0.01) & 0.66 (0.00) & 0.667 \\
Read & 65,536 & 1.00 (0.00) & 0.74 (0.08) & 0.747 \\
Write/read & 256 & 0.79 (0.10) & 0.55 (0.01) & 0.698 \\
Write/read & 4,096 & 0.87 (0.01) & 0.53 (0.00) & 0.610 \\
Write/read & 65,536 & 0.97 (0.11) & 0.58 (0.02) & 0.600 \\
\hline
\end{tabular}
\end{center}


These timings compare two complete implementations. Managed storage,
bounds checks, wrapper code, and compiler optimization can all affect the
result. They do not isolate the cost of a type parameter, measure allocation
cost, or predict performance for other element types. They are a new paired
run, separate from the historical standard-library comparison in
Section~\ref{sec:benchmarks}. The raw samples, summaries, compiler version,
machine details, source hashes, and plots are saved under
\path{example/poly_array/benchmark/results/}.
